\scanpage{29}
\mspage{- 214 -}\hfill n\textsuperscript{o} 391

son ordre est $h-1=17$. Elle est orthogonale à toutes les racines simples sauf à $\alpha_1$ et l'on a $(\tilde\alpha\mid\alpha_1)=1$.

(V) Ce qui a été vu au n\textsuperscript{o}~10 montre que les \uline{coracines} de $R'$ s'obtiennent encore en remplaçant $\varepsilon_i$ par $\varepsilon_i^*$ dans les expressions des racines de $R'$, et $R'{}^*$ est donc isomorphe à $R'$. Toutes les racines ayant même longueur, la \uline{forme bilinéaire canonique} pour $R$ est $(x\mid y)/2h=(x\mid y)/36$ et la longueur des racines est $1/3\sqrt{2}$, (\S~8, n\textsuperscript{o}~5, \uline{Remarque}). Le nombre des racines non orthogonales à une racine donnée est $4h-6=66$ et le nombre $\gamma$ est égal à $h^2=2^2.3^4$.

(VI) Le calcul des poids fondamentaux, formant la base duale de $(\alpha_i^*)$, $(1\leq i\leq7)$, dans $V'$, donne :

\[
\begin{aligned}
\omega_1&=\varepsilon_8-\varepsilon_7=2\alpha_1+2\alpha_2+3\alpha_3+4\alpha_4+3\alpha_5+2\alpha_6+\alpha_7=\tilde\alpha\\
\omega_2&=\frac{1}{2}(\varepsilon_1+\varepsilon_2+\varepsilon_3+\varepsilon_4+\varepsilon_5+\varepsilon_6-2\varepsilon_7+2\varepsilon_8)=\\
&=\frac{1}{2}(4\alpha_1+7\alpha_2+8\alpha_3+12\alpha_4+9\alpha_5+6\alpha_6+3\alpha_7)\\
\omega_3&=\frac{1}{2}(-\varepsilon_1+\varepsilon_2+\varepsilon_3+\varepsilon_4+\varepsilon_5+\varepsilon_6-3\varepsilon_7+3\varepsilon_8)=\\
&=3\alpha_1+4\alpha_2+6\alpha_3+8\alpha_4+6\alpha_5+4\alpha_6+2\alpha_7\\
\omega_4&=\varepsilon_3+\varepsilon_4+\varepsilon_5+\varepsilon_6+2(\varepsilon_8-\varepsilon_7)=\\
&=4\alpha_1+6\alpha_2+8\alpha_3+12\alpha_4+9\alpha_5+6\alpha_6+3\alpha_7\\
\omega_5&=\frac{1}{2}(2\varepsilon_4+2\varepsilon_5+2\varepsilon_6+3(\varepsilon_8-\varepsilon_7))=\\
&=\frac{1}{2}(6\alpha_1+9\alpha_2+12\alpha_3+18\alpha_4+15\alpha_5+10\alpha_6+5\alpha_7)\\
\omega_6&=\varepsilon_5+\varepsilon_6-\varepsilon_7+\varepsilon_8=2\alpha_1+3\alpha_2+4\alpha_3+6\alpha_4+5\alpha_5+4\alpha_6+2\alpha_7\\
\omega_7&=\varepsilon_6+\frac{1}{2}(\varepsilon_8-\varepsilon_7)=\frac{1}{2}(2\alpha_1+3\alpha_2+4\alpha_3+6\alpha_4+5\alpha_5+4\alpha_6+3\alpha_7).
\end{aligned}
\]

% archive sheet number: 28
